


\begin{comment}


\begin{figure}[t]
    \centering
    \includegraphics[width=0.99\linewidth]{figures/overview-figure.png}
    % \vspace{-.5pc}
    \caption{The figure captioning task can be tackled as a text-summarization task. The figure's caption can be generated by summarizing the paragraphs mentioning the figure. (Paper:~\cite{doulaty2015data}.)}
    \label{fig:overview}
    % \vspace{-1pc}
\end{figure}

\end{comment}

\begin{figure*}[t]
    \centering
    \includegraphics[width=0.99\textwidth]{figures/overview-figure-teaser-v2.png}
    \vspace{-3mm}
    \caption{Figure captioning can be addressed as a text-summarization task. The figure’s caption can be generated by summarizing the paragraphs mentioning the figure. The caption is generated by the model Pegasus$_{P+O+B}$. The example shown in this figure is extracted from the paper~\cite{doulaty2015data}.}
   \label{fig:overview}
   \vspace{-2mm}
\end{figure*}

In scientific documents, effective figure captions help readers understand complex figures like bar charts, line charts, or pie charts.
These captions describe the images and often include necessary context from the document's full text~\cite{durbin2004effective}.
Unfortunately, even published papers often have poorly-written captions.
As per our analysis (\Cref{sec:quality-annotation}), around 53.88\% of line charts' captions in arXiv \texttt{cs.CL} papers are found to be unhelpful for NLP readers.
Automatic caption generation could aid paper writers by providing good starting captions that can be refined for better quality.

%Prior work often treated figure caption generation as a \textit{vision-to-language} task, \ie, generating a caption primarily based on the image. 
Previous research typically approached figure caption generation as a \textit{vision-to-language} task, \ie, creating captions based on the image.
%For example, \citet{hsu-etal-2021-scicap-generating} took an end-to-end approach based on CNN+RNN structures, which %directly 
For instance, \citet{hsu-etal-2021-scicap-generating} used an end-to-end approach with CNN+RNN structures, which extracted feature representation from the image and converted it into caption text. 
%extracts feature representation from the image, and then decoded them into a descriptive text sequence. 
%\citet{qian2021generating} tackled the issue from an image-understanding perspective, first extracting semantic information from the image and using a template for generation.
\citet{qian2021generating} took a slightly different approach: first understanding what is in the image, pulling out key information, and then using a preset template to create the caption.
However, although achieving some success in synthetic data~\cite{kahou2017figureqa,kafle2018dvqa,chen2019figure,zhu2021autochart}, these approaches often struggled to caption real-world figures.
For example, \citet{hsu-etal-2021-scicap-generating}'s end-to-end approach, trained and tested using arXiv figures, achieved a BLEU-4 score of only 2.91.

In this paper, we argue that figure captioning in scientific documents can be more effectively tackled as a \textbf{text-summarization task}: 
The caption can be generated by \textbf{summarizing the paragraphs mentioning the figure} (as shown in Figure~\ref{fig:overview}.)
Scientific figures typically come with extensive text in the scientific document that can aid caption generation.
Our analysis (Section~\ref{sec:motivating-analysis}) shows that, in arXiv, over 75\% of words in figure captions can be aligned with the words in the paragraphs referencing those figures, which motivates our approach.
%inspired us to use language technologies to produce captions from these paragraphs.
%Figure~\ref{fig:overview} overviews our approach, extracting sentences referencing the target figure and then summarizing them into a concise caption.
The automatic evaluation shows that summarizing paragraphs referencing the figures results in better captions than prior vision-based methods.
%Even simply using sentences that mention the figures-- without any changes-- as captions yielded higher ROUGE scores.
In a human evaluation by external domain experts, our best-performing model's captions were preferred over the original captions 46.67\% of the time.

%We further investigated two major challenges: {\em (i)} low-quality author-written captions and {\em (ii)} the lack of clear standards for good captions.
We further conducted an in-depth investigation focused on two key challenges: {\em (i)} the common presence of low-quality author-written captions and {\em (ii)} the lack of clear standards for good captions.
Surprisingly, 53.88\% of the author-written captions in our sample was deemed unhelpful.
This has implications for the design of future captioning systems, underscoring the influence of data quality on captioning performance.


%-------------------------------------------------------------------------
\begin{comment}


The main contributions of this work are summarized as follows.

\begin{itemize}[leftmargin=*,itemsep=0mm]

    \item Revolutionizing figure captioning in scholarly articles by showing the effectiveness of text summarization over traditional vision approaches. A novel combination of existing data mining tools brings a fresh perspective to the task.

    \item A quantitative measurement by domain experts found that over half of the figure captions in arXiv papers are unhelpful, highlighting the real-world data quality challenge in the figure captioning task (Section~\ref{sec:analysis-challenge-1}).

    \item A data-driven analysis clarifies the relationships between text, captions, and figures, offering benefits for the future development of vision-and-language tasks related to scientific figures and articles (Section~\ref{sec:motivating-analysis}).
    
    \item An examination of the relationships between factors (such as the presence of high-level takeaway messages) and reader-perceived caption quality reveals what constitutes a good caption (Section~\ref{sec:quality-annotation}). 
This lays the foundation for optimizing figure captioning models.

\end{itemize}

    
\end{comment}



%------------------- dead kitten --------------

\begin{comment}



%----------------- KDD draft below ------------------
\bigskip

Effective figure captions are vital for clear comprehension of scientific figures.
These captions help readers understand the intended message by describing the images and often include necessary context from the document's full text~\cite{durbin2004effective}.
However, poor caption writing remains a widespread issue in scientific articles despite its importance.
%Clear and effective figure captions are critical for effective communication of scientific figures. These captions play a crucial role in providing context and imparting the intended message of the visual representation to the reader. Despite their significance, poor caption writing continues to be a prevalent challenge in the scientific community~\cite{durbin2004effective}. 
Our analysis (\Cref{sec:quality-annotation}) showed that approximately 53.88\% of line charts in arXiv \texttt{cs.CL} papers are unhelpful to NLP readers.
Thus, our key motivation is to provide better captions to readers when the original caption is absent or poorly written.

While several prior works attempted to generate captions for figures automatically, 
they often treated the problem as a vision-to-language task, \ie, generating a caption primarily based on the image. 
For example, \citet{hsu-etal-2021-scicap-generating} took an end-to-end approach based on CNN+RNN structures, which %directly 
extracts feature representation from the image, and then decoded them into a descriptive text sequence. 
\citet{qian2021generating} tackled the issue from an image-understanding perspective, first extracting semantic information from the image and using a template for generation.
However, these approaches struggled to caption real-world figures. 
The end-to-end model trained on \oldDataset achieved a BLEU-4 score of only 2.91; many other works solely experimented with synthetic figures~\cite{kahou2017figureqa,kafle2018dvqa,chen2019figure,zhu2021autochart}.
Furthermore, while these approaches make sense when no textual information is provided, scientific figures typically have much associated textual information that can be utilized to enhance the caption-generation process.
To the best of our knowledge, no existing figure-caption datasets explicitly contain the figures' accompanying documents~\cite{ImageCLEFmedConceptOverview2021,hsu-etal-2021-scicap-generating,10.1145/3341162.3345601} because prior work treated this problem as a vision-to-language~\cite{9973197,ImageCLEFmedConceptOverview2021,hsu-etal-2021-scicap-generating,qian2021generating,10.1145/3341162.3345601,kanthara2022chart} or image-understanding task~\cite{kahou2017figureqa,kafle2018dvqa,Methani_2020_WACV}. 
%\kenneth{TODO Kenneth: Move the rebuttal texts here, saying some datasets do not even have the article available.}\cy{done}%



%Furthermore, while these approaches are suitable for cases where there is no accompanying textual information, scientific figures typically have abundant associated textual information that can be utilized to enhance the caption generation process.
%while these approaches make sense when no textual information is provided, they did not exploit the rich textual information often available in the accompanying document to generate good captions.
%\sukim{what about mentioning use-cases and limitations of them? It makes sense to solve this as an img2txt problem when no textual information is provided. But when it comes to webpages or scientific figures, they usually have rich textual information that we could leverage, and there is no reason not to exploit this information to generate good captions. In addition, when they use OCR, they often struggle with incorrect recognition. We could brought the preliminary Mturk study showing that TrOCR performs the worst.}

In this paper, we argue that figure captioning in articles can be effectively tackled as a \textbf{text-summarization task}: The caption can be generated by \textbf{summarizing the paragraphs mentioning the figure}.
Our preliminary analysis, detailed in Section~\ref{sec:motivating-analysis}, reveals that over 75\% of the tokens in figure captions can be aligned with the corresponding paragraphs that mention the figures.
This observation motivates the use of language technologies for caption generation based on these paragraphs.
%This approach stems from a series of our motivating studies (Section~\ref{sec:motivating-analysis}) showing that over 75\% of the tokens in figure captions can be aligned with the corresponding paragraphs that mention the figure, suggesting the effectiveness of using language technologies to solve task.
%
Figure~\ref{fig:overview} shows our approach that extracts sentences referencing the target figure, then summarizes them into a concise caption.
The automatic evaluation of our experiments with 
%that automatically evaluated articles in 
the \oldDataset dataset showed that summarizing the figure-mentioning paragraphs produced significantly better captions than traditional vision-based approaches.
%, per automatic evaluation.
Even simply using the figure-mentioning sentences without any modifications as the captions earned better ROUGE scores.
In a human evaluation conducted by external domain experts, our top-performing model's captions were chosen over the original captions 46.67\% of the time.

%To move the needle in figure captioning and inform the development of future caption generation systems, we further conducted in-depth investigations into the two core challenges: 
%{\em (i)} low-quality author-written captions and 
%{\em (ii)} the absence of a standard for good captions.
To truly move the needle in figure captioning, we further investigated the two core real-world challenges of the problem: 
{\em (i)} low-quality author-written captions and 
{\em (ii)} the absence of a standard for good captions.
We manually rated the usefulness of 438 captions 
sampled from \oldDataset and identified that 53.88\% of the captions were considered unhelpful. 
For the figures with original captions rated as unhelpful, our best-performing model can often generate better-quality captions.
Our analysis also found that good captions tend to include the high-level takeaway of the figure.
These findings have implications for the future development and design of caption generation systems, providing deeper insights into the impact of data quality on figure captioning performance.



\kenneth{KENNETH IS EDITING HERE}


\bigskip
%------------------------- old draft ------------------

%What is caption and why they're important
%Scientists use charts in papers to show experimental results; companies use charts in business reports to show financial status; journalists use charts in news articles to show trends and statistics.~
Given the abundance of documents on the Internet,
figures play an important role in enriching technical articles. 
By reading the illustrative figures from online resources such as arXiv or newspaper websites,
we understand scientific results better,
interpret business reports from companies deeper,
and remember trends from news articles easier.
% Scientists use figures to show experimental results, 
% companies use figures in business reports to demonstrate financial status, 
% and journalists use figures in news articles to display trends and statistics. 
%\kenneth{Can use one more or one better example here where *a chart is in a document*.}\ed{government use chart in the official documents to show the trends and disease statistics}.
%Visuals such as figures, graphics, and diagrams help readers understand the content of a full-length document effectively.
%The captions of these visuals help readers understand the intended message by describing the images and often include necessary context from the document's full text~\cite{durbin2004effective}. 


% Kenneth: I moved the following to Method section.
%\ed{\citet{hartley2003single} and \citet{gelman2002let} both argue
%that increasing the caption lengths to explain more fully what the data actually show can
%increase the readers' comprehension of both tables and figures}
%\kenneth{TODO Ed: Find relevant citations about the role of captions, e.g., author guideline for journals}

% Gromit: I did some re-arrangement of sentences here.
%Several prior works attempted to generate captions for figures automatically. 
%A key motivation is to support readers, especially when the original caption is absent or poorly written.
However, captions in current academic papers seem to be ineffective for readers. 
Our analysis (\Cref{sec:quality-annotation}) showed that approximately 50\% of line charts in arXiv \texttt{cs.CL} papers are unhelpful to NLP readers.\kenneth{This paper has too many ``NLP'' in it. Maybe replace it with something else.} \ed{maybe use domain experts} \sukim{could be simply readers in this case}
Thus, our key motivation is to provide better captions to readers when the original caption is absent or poorly written.
While several prior works attempted to generate captions for figures automatically, 
they often treated the problem as a vision-to-language task, \ie, generating a caption primarily based on the image. 
For example, \citet{hsu-etal-2021-scicap-generating} took an end-to-end approach based on CNN+RNN structures, which directly extract feature representation from the image, and then decoded them into a descriptive text sequence. 
\citet{qian2021generating} tackled the issue from an image-understanding perspective, first extracting semantic information from the image and using a template for generation.
However, these approaches often struggled to caption real-world figures. 
The end-to-end model trained on \oldDataset achieved a BLEU-4 score of only 2.91; many other works solely experimented with synthetic figures~\cite{kahou2017figureqa,kafle2018dvqa,chen2019figure,zhu2021autochart}.
\sukim{what about mentioning use-cases and limitations of them? It makes sense to solve this as an img2txt problem when no textual information is provided. But when it comes to webpages or scientific figures, they usually have rich textual information that we could leverage, and there is no reason not to exploit this information to generate good captions. In addition, when they use OCR, they often struggle with incorrect recognition. We could brought the preliminary Mturk study showing that TrOCR performs the worst.}

In this paper, we argue that figure captioning in articles can be tackled as a \textbf{text-summarization task}: The caption can be generated by \textbf{summarizing the paragraphs mentioning the figure}. (See Figure~\ref{fig:overview}.)
%
Experiments that automatically evaluated articles in the \oldDataset dataset showed that summarizing the figure-mentioning paragraphs produced significantly better captions than \textit{vision-to-language} approaches. 
Even simply using the figure-mentioning sentences without any modifications as the captions earned better ROUGE scores.

Three NLP researchers (who are not co-authors of this paper) were recruited to rank captions. 
They preferred our best-performing model's caption over the original caption 46.67\% of the time. 
We also investigated the relationships between the captions and the paragraphs that mention them. 
Around 76.68\% of tokens in an average figure caption can be aligned to words in the figure-mentioning paragraphs; 
early evidence suggested that the remaining 23.32\% might refer to visual information available only in the image.

To gain deeper insight into the problem, we manually rated the usefulness of 438 captions 
sampled from \oldDataset and found that 53.88\% of the captions were considered unhelpful to readers. 
In cases where the original caption was unhelpful, 
the NLP researchers preferred machine-generated captions 22 out of 35 times.
This work presents a novel and practical approach to automatically generating figure captions. 
Our findings can inform future research in developing new technologies to help caption readers and writers,
\eg, new tools that present figures with comprehensive captions to help readers quickly catch important messages from scientific papers.


\kenneth{Need to list three contributions explicitly: (1) New perspective to the task (quote KDD CFP), (2) make the relations between text, caption, and figures more clear, and (3) show that author-written captions can often be bad!}


% \kenneth{TODO CY: UPDATE NUMBERS}

%\kenneth{TODO: Need a few sentences to wrap up the implications.}



Three NLP researchers (who are not co-authors of this paper) were recruited to rank captions. 
They preferred our best-performing model's caption over the original caption 46.67\% of the time. 



%----------------- before editor editing -------------
In this paper, we argue that the figure captioning task can be tackled as a \textbf{text summarization task}.
When a figure is part of an article, the figure's caption can be generated by \textbf{summarizing the paragraphs mentioning the figure} (See Figure~\ref{fig:overview}.)
%Summarizing the paragraph-- or even the sentence-- that mentions the figure is more effective in producing captions for scientific articles.
%Our experiments showed that the summarization approach drastically outperformed all other baselines in both automatic and human evaluations.
%Our extensive experiments 
Our experiments on the \oldDataset dataset showed that summarizing the figure-mentioning paragraphs produces significantly better captions than vision-to-language approaches, per automatic evaluation.
Even simply using the figure-mentioning sentence-- without any modifications-- as the captions reached better ROUGE scores.
%For human evaluation, we recruited 
Three NLP researchers (who are not co-authors of this paper) were recruited to rank captions, where
%On average, 
our best-performing model's caption was preferred over the original caption 46.67\% of the time.
% \kenneth{\#(long-gen wins origin even when long-gen is not the best)/\#all}
We also investigated the relations between the captions and the paragraphs that mention them.
Around 76\% of tokens in an average figure caption can be aligned to words in the figure-mentioning paragraphs; early evidence suggested that the remaining 24\% might refer to visual information that is only available in the image.

%all wins analysis 90 / 42
%helpful wins analysis 55 / 24
%unhelpful wins analysis 35 / 18

%%% machine-generated captions is preferred
% 70 / 105 (all annotation)
% 22 / 35 (averaging)

%Finally, 
To gain deeper insight into the problem space, 
we manually rated the usefulness of 438 captions sampled from \oldDataset.
%randomly extracted from arXiv papers. %\ed{all arxiV papers or paper in SciCap?}
As a result, 53.88\% of the captions were considered not helpful to readers.
We also found that the NLP researchers often preferred machine-generated captions, \ie, 22 out of 35 times, for these cases where the original caption is unhelpful.

This work presents a novel and practical approach to generating figure captions. 
Our findings can inform future research to develop new technologies to help caption readers and writers.



%----------------- old intro below -------------
\bigskip
%----------------- old intro below -------------


For human evaluation, we recruited three NLP researchers (who are not co-authors of this paper) to rank captions.
On average, our best-performing model's caption was preferred over the original caption XX\% of the time.\kenneth{\#(long-gen wins origin)/\#all}
When the original caption is not helpful, the generated captions were even more preferable, \ie, YY\% of the time.\kenneth{\#(long-gen or gen wins origin)/\#all *in the NOT HELPFUL batch*}\kenneth{TODO CY: update the number (Are the numbers legit?)}





%what's our cool contribution
%gromit: for me the coolest contribution is the way we evaluate it. The evaluation method should be highlighted here.
Figure captions can be approximated by summarizing or paraphrasing mentions/paragraphs, 
%gromit: the below sentence seems too aggressive, consider remove it
raising questions about the role of figure images in caption generation.
Further analysis demonstrated that 75\% of tokens in an average figure caption can be aligned to words in the figure-mentioning paragraph.

%cherry on the top


We also revealed that many low-quality figure captions exist in arXiv paper.
%A side bonus could be some: 
In the case where human-written captions are bad, we can generate better captions using mentions.
\kenneth{Is there additional information in the sentence mentioning the figure? the first and second sentence? and so on.}

%The caption of a figure can be reliably generated by \textbf{summarizing the paragraphs mentioning the figure}.





Several prior works attempted to generate captions for figures automatically. 
One key motivation is to support readers, especially when the original caption is absent or poorly written.
Our analysis (Section~\ref{}) showed that approximately 50\% of line charts in arXiv {\tt cs.CL} papers are not helpful to NLP readers.~\kenneth{UPDATE SECTION REF and update numbers}


Prior work attempted to automatically generate captions for figures or charts, motivated by the fact that many low-quality captions exist even in published articles.
Our analysis (Section~\ref{}) showed that over 50\% of line charts in arXiv {\tt cs.CL} papers are considered not helpful.\kenneth{UPDATE SECTION REF and update numbers}
Generating chart captions has been treated as a \textit{vision-to-language} task, \ie, the model generates a caption primarily based on the image.
For example, \kenneth{DESCRIBE OUR ONE WORK}
Some other approaches first extract pieces of information from the image and put together all information using a template\kenneth{CITE JOEL'S WORK and Lee's earlier work}.
However, such approaches struggled with real-world figures.
The end-to-end model trained on \oldDataset only achieved a BLEU-4 score of 2.91, which is far from practical use.
Many prior works restrained themselves from experimenting with real-world charts and solely focused on synthetic data. 

%In this paper, we argue that chart captioning can be viewed as a \textbf{text summarization task} when the chart is a part of an article or report.
%\textbf{Summarizing the article's paragraphs that mention the chart can generate a high-quality caption for the chart.}




Several prior works attempted to automatically generate captions for figures or charts, motivated by the fact that many low-quality captions exist even in published articles~\cite{hsu-etal-2021-scicap-generating}.
Our analysis showed that 50\% of line charts in arXiv {\tt cs.CL} papers are considered not helpful or worse (Section~\ref{}.)\kenneth{UPDATE SECTION REF and update numbers}




%what is known
Generating chart captions has been treated as a vision-to-language task, \ie, the model produces a caption primarily based on the image.
However, when working with real-world figures in scientific articles, such models do not produce good-enough captions that people could use~
\cite{hsu-etal-2021-scicap-generating}.
\kenneth{Also mention many prior work used synthetic data.}

%why it's not enough

This paper argues that the figure captioning task in scientific articles can be viewed as a \textbf{text summarization task} instead of a vision-to-language task.

\end{comment}